\subsection{典范双线性形式}

我们已经看到（第 1 号，命题 3），对称双线性形式

\[
(x, y) \mapsto \mathrm{B} _ {\mathrm{R}} (x, y) = \sum_ {\alpha \in \mathbb {R}} \langle \alpha^ {\prime}, x \rangle \langle \alpha^ {\prime}, y \rangle
\]

它在 V 上非退化且在 A(R) 下不变。交换 R 与 \(R^{\prime}\) 的作用，便得到，对称双线性形式 \((x^{*}, y^{*}) \mapsto \mathrm{B}_{\mathbb{R}^{\prime}}(x^{*}, y^{*}) = \sum_{\alpha \in \mathbb{R}^{\prime}} \langle \alpha, x^{*} \rangle \langle \alpha, y^{*} \rangle\) 在 \(V^{*}\) 上非退化且在 A(R) 下不变。

\(B_{R^{\prime}}\)（分别为 \(B_{R}\)）在 V（分别为 \(V^{*}\)）上的逆形式称为 V（分别为 \(V^{*}\)）上的典范双线性形式，记为 \(\Phi_{R}\)（分别为 \(\Phi_{R^{\prime}}\)）。它非退化且在 A(R) 下不变。令 \(\sigma\) 为从 V 到 \(V^{*}\) 且由 \(B_{R^{\prime}}\) 定义的同构。于是，对于 \(x \in V\) 和 \(y \in V\)：

\[
\Phi_ {R} (x, y) = B _ {R ^ {\prime}} (\sigma (x), \sigma (y)) = \sum_ {\alpha \in R} \langle \alpha , \sigma (x) \rangle \langle \alpha , \sigma (y) \rangle .
\]

但是 \(\langle \alpha, \sigma(x) \rangle = B_{R^{\prime}}(\sigma(\alpha), \sigma(x)) = \Phi_R(\alpha, x)\)。因此，

\[
\Phi_ {\mathsf {R}} (x, y) = \sum_ {\alpha \in \mathsf {R}} \Phi_ {\mathsf {R}} (\alpha , x) \Phi (\alpha , y).\tag{16}
\]

根据第 2 号命题 7，\(\Phi_{R}\) 是满足恒等式 (16) 且在 W(R) 下不变的唯一非零对称双线性形式。

对于 \(\beta \in \mathbb{R}\)，由 (16) 得

\[
\varPhi_ {\mathrm{R}} (\beta , \beta) = \sum_ {\alpha \in \mathrm{R}} \varPhi_ {\mathrm{R}} (\alpha , \beta) ^ {2} = \frac {1}{4} \varPhi_ {\mathrm{R}} (\beta , \beta) ^ {2} \sum_ {\alpha \in \mathrm{R}} n (\alpha , \beta) ^ {2},
\]

故

\[
4 \Phi_ {R} (\beta , \beta) ^ {- 1} = \sum_ {\alpha \in R} n (\alpha , \beta) ^ {2}.\tag{17}
\]

此外，根据第 1 号引理 2，对于 \(x, y \in V\)，有：

\[
\begin{array}{r l} & {\mathrm{B} _ {\mathrm{R}} (x, y) = \sum_ {\alpha \in \mathrm{R}} \varPhi_ {\mathrm{R}} \left(\frac {2 \alpha}{\varPhi_ {\mathrm{R}} (\alpha , \alpha)}, x\right) \left(\frac {2 \alpha}{\varPhi_ {\mathrm{R}} (\alpha , \alpha)}, y\right)} \\ & {\qquad = 4 \sum_ {\alpha \in \mathrm{R}} \varPhi_ {\mathrm{R}} (\alpha , x) \varPhi_ {\mathrm{R}} (\alpha , y) \varPhi_ {\mathrm{R}} (\alpha , \alpha) ^ {- 2}.} \end{array}
\]

由此，若 \(R\) 不可约，则存在常数 \(\gamma(R) > 0\)，使得

\[
\sum_ {\alpha \in R} \Phi_ {R} (\alpha , x) \Phi_ {R} (\alpha , y) \Phi_ {R} (\alpha , \alpha) ^ {- 2} = \gamma (R) \Phi_ {R} (x, y).\tag{18}
\]

根据 \(\gamma (\mathbb{R})\) 的定义，有 \(B_{\mathbb{R}}(x,y) = 4\gamma (\mathbb{R})\Phi_{\mathbb{R}}(x,y)\)，故

\[
\Phi_ {\mathrm{R} ^ {\prime}} \left(x ^ {*}, y ^ {*}\right) = (4 \gamma (\mathrm{R})) ^ {- 1} \mathrm{B} _ {\mathrm{R} ^ {\prime}} \left(x ^ {*}, y ^ {*}\right)
\]

对于 \(x^{*},y^{*}\in V^{*}\)。这证明了 \(\gamma (\mathsf{R}) = \gamma (\mathsf{R}^{\prime})\)。另一方面，对于 \(\beta \in \mathbb{R}\)，

\[
\begin{array}{c} \Phi_ {\mathrm{R} ^ {\prime}} (\beta^ {\prime}, \beta^ {\prime}) = (4 \gamma (\mathrm{R})) ^ {- 1} \sum_ {\alpha \in \mathfrak {R}} \langle \beta^ {\prime}, \alpha \rangle^ {2} \\ = \gamma (\mathrm{R}) ^ {- 1} \sum_ {\alpha \in \mathbb {R}} \frac {\Phi_ {\mathrm{R}} (\alpha , \beta) ^ {2}}{\Phi_ {\mathrm{R}} (\beta , \beta) ^ {2}} \end{array}
\]

所以，根据 (16)，

\[
\Phi_ {\mathrm{R} ^ {\prime}} (\beta^ {\prime}, \beta^ {\prime}) = \gamma (\mathrm{R}) ^ {- 1} \Phi_ {\mathrm{R}} (\beta , \beta) ^ {- 2} \Phi_ {\mathrm{R}} (\beta , \beta)
\]

或最后

\[
\Phi_ {R} (\beta , \beta) \Phi_ {R ^ {\prime}} (\beta^ {\prime}, \beta^ {\prime}) = \gamma (R) ^ {- 1}.\tag{19}
\]

此外，若 \(\mathbb{R}\) 的所有根相对于 \(\lambda\) 都具有相同长度（长度由 \(\varPhi_{\mathbb{R}}\) 给出），则 (16) 和 (18) 表明

\[
\gamma (R) = \lambda^ {- 4}.\tag{20}
\]

此外，若 \(h\) 是 \(W\) 的 Coxeter 数，则第五章第 6 节第 2 号定理 1 的推论表明：

\[
h \Phi_ {R} (x, x) = \sum_ {\alpha \in R} \left(x, \frac {\alpha}{\lambda}\right) ^ {2} \text {   for   all   } x \in V.
\]

与 (16) 比较，得到

\[
\lambda = h ^ {- 1 / 2} \text { or } \gamma (\mathrm{R}) = h ^ {2}.\tag{21}
\]

最后，公式 (19) 表明，\(R^{\prime}\) 的根具有长度 \(\lambda\)（相对于 \(\Phi_{R^{\prime}}\)）。

